\chapter{Messergebnisse}
\label{ch:messergebnisse}

Dieses Kapitel vergleicht die Messergebnisse der drei untersuchten Konfigurationsprofile Minimal, Logging und IoT-Stack. Jede zeitbezogene Messreihe beinhaltet dabei 1000 ausgewertete Messwerte nach zehn verworfenen Aufw{\"a}rmdurchl{\"a}ufen. Die DWT-basierten Ergebnisse werden in CPU-Zyklen bei einer Taktfrequenz von 64 MHz angegeben. Ein CPU-Zyklus entspricht damit 15,625 ns. Der Timer-Benchmark benutzt den RTC-Systemtimer mit einer Frequenz von 32768 Hz.

Vor den eigentlichen Benchmarks wurde in jedem Profil ein Selbsttest mit einer vorgesehenen Wartezeit von 100 \textmu{}s durchgef{\"u}hrt. Die Mittelwerte lagen zwischen 6432,01 und 6435,96 CPU-Zyklen bzw. zwischen 100,50 und 100,56 \textmu{}s. Damit wurde das in Kapitel 5 festgelegte Abnahmekriterium in allen drei Profilen erf{\"u}llt.

Tabelle~\ref{tab:selbsttest-100us} fasst den im Text dokumentierten Wertebereich der drei Selbsttests zusammen. Da die Ausgangsdaten an dieser Stelle nur den kleinsten und gr{\"o}{\ss}ten Profilmittelwert ausweisen, werden keine nicht belegten Einzelwerte je Profil erg{\"a}nzt.

\begin{table}[H]
\centering
\caption{Ergebnisse des Selbsttests mit einer Sollzeit von 100 \textmu{}s}
\label{tab:selbsttest-100us}
\begin{tabular}{lrr}
\toprule
Auswertung & Mittelwert in CPU-Zyklen & Mittelwert in \textmu{}s \\
\midrule
Kleinster Profilmittelwert & 6432,01 & 100,50 \\
Gr{\"o}{\ss}ter Profilmittelwert & 6435,96 & 100,56 \\
\bottomrule
\end{tabular}
\end{table}

Die in Tabelle~\ref{tab:selbsttest-100us} dargestellte Spannweite betr{\"a}gt 3,95 CPU-Zyklen beziehungsweise 0,06 \textmu{}s. Beide Grenzwerte liegen damit nahe an der vorgegebenen Sollzeit und best{\"a}tigen das definierte Abnahmekriterium.

\section{Kontextwechsel}
F{\"u}r die Kontextwechselmessung wurden zwei Varianten untersucht. Tabelle~\ref{tab:kontextwechsel-profile} vergleicht ihre mittleren Latenzen f{\"u}r alle drei Konfigurationsprofile. Untersucht wurden die Aktivierung eines h{\"o}herprioren Threads durch die Freigabe eines Semaphors und die kooperative {\"U}bergabe zwischen zwei gleichprioren Threads durch \lstinline[style=codeinline]|k_yield()|.

\begin{longtable}{p{0.313\linewidth}p{0.313\linewidth}p{0.313\linewidth}}
\caption{Kontextwechsel-Latenzen der drei Konfigurationsprofile}\label{tab:kontextwechsel-profile}\\
\toprule
\textbf{Profil} & \textbf{Semaphore-Weckung} & \textbf{Kooperatives Yield} \\
\midrule
Minimal & 407,86 Zyklen & 228,00 Zyklen \\
Logging & 439,99 Zyklen & 240,00 Zyklen \\
IoT-Stack & 428,96 Zyklen & 240,99 Zyklen \\
\bottomrule
\end{longtable}

Tabelle~\ref{tab:kontextwechsel-profile} zeigt f{\"u}r beide Varianten einen Anstieg gegen{\"u}ber dem Minimalprofil, wobei die Semaphore-Weckung in allen Profilen langsamer als das kooperative Yield ist.

Die Semaphore-Weckung braucht im Minimalprofil durchschnittlich 407,86 Zyklen bzw. 6,37 \textmu{}s. Im Logging-Profil steigt der Mittelwert auf 439,99 Zyklen bzw. 6,87 \textmu{}s. Dies entspricht einer Zunahme um 7,88 \%. Das IoT-Profil erreicht dabei 428,96 Zyklen bzw. 6,70 \textmu{}s und liegt damit 5,17 \% {\"u}ber dem Minimalprofil.

Beim kooperativen Yield betr{\"a}gt der Mittelwert im Minimalprofil 228,00 Zyklen beziehungsweise 3,56 \textmu{}s. Logging braucht 240,00 Zyklen und der IoT-Stack 240,99 Zyklen. Im Vergleich zum Minimalprofil entsprechen diese Werte Zunahmen um 5,26 \% bzw. 5,70 \%.

\section{Interrupt-Latenz}
Die Interrupt-Latenz wurde vom Setzen der NVIC-Anforderung bis zum ersten Zeitstempel im C-Handler gemessen. Tabelle~\ref{tab:interrupt-profile} stellt Minimum, Maximum, Mittelwert und Standardabweichung der drei Profile gegen{\"u}ber.

\begin{longtable}{p{0.188\linewidth}p{0.188\linewidth}p{0.188\linewidth}p{0.188\linewidth}p{0.188\linewidth}}
\caption{Interrupt-Latenzen der drei Konfigurationsprofile}\label{tab:interrupt-profile}\\
\toprule
\textbf{Profil} & \textbf{Minimum} & \textbf{Maximum} & \textbf{Mittelwert} & \textbf{Standardabweichung} \\
\midrule
Minimal & 41 & 47 & 46,99 & 0,19 \\
Logging & 38 & 38 & 38,00 & 0,00 \\
IoT-Stack & 37 & 45 & 44,94 & 0,56 \\
\bottomrule
\end{longtable}
Wie Tabelle~\ref{tab:interrupt-profile} zeigt, besitzt das Logging-Profil den niedrigsten und zugleich konstanten Messwert; das Minimal- und IoT-Profil weisen dagegen eine kleine Streuung auf.

Das Minimalprofil erreicht im Mittel 46,99 Zyklen bzw. 0,734 \textmu{}s. Beim Logging-Profil wurden dabei 38,00 Zyklen bzw. 0,594 \textmu{}s und beim IoT-Profil 44,94 Zyklen bzw. 0,702 \textmu{}s gemessen. Logging liegt damit 19,13 \% und der IoT-Stack 4,36 \% unter dem Minimalprofil. Die Ursachen dieser nicht monotonen Entwicklung werden in Kapitel 8 diskutiert.

\section{Timer-Jitter}
F{\"u}r eine angeforderte Periode von 10 ms ergibt sich bei einer RTC-Frequenz von 32768 Hz ein Sollwert von 328 Systemtimerzyklen. Ein RTC-Zyklus entspricht dabei ungef{\"a}hr 30,52 \textmu{}s. Tabelle~\ref{tab:timer-perioden-profile} zeigt die gemessenen Perioden aller drei Profile.

\begin{table}[H]
\centering
\caption{Timer-Perioden der drei Konfigurationsprofile}
\label{tab:timer-perioden-profile}
\begin{tabular}{lrrr}
\toprule
Profil & Mittelwert in RTC-Zyklen & Standardabweichung & Periodendauer in ms \\
\midrule
Minimal & 328,00 & 0,00 & 10,0098 \\
Logging & 328,00 & 0,00 & 10,0098 \\
IoT-Stack & 328,00 & 0,00 & 10,0098 \\
\bottomrule
\end{tabular}
\end{table}

Wie Tabelle~\ref{tab:timer-perioden-profile} verdeutlicht, wurden in allen drei Profilen f{\"u}r mehrere 1000 Perioden jeweils genau 328 RTC-Zyklen gemessen. Innerhalb der verf{\"u}gbaren RTC-Aufl{\"o}sung war damit kein Jitter und kein Unterschied zwischen den Konfigurationsprofilen sichtbar.

\section{Heap-Allokationsgeschwindigkeit}
Es wurden die Anforderung und die anschlie{\ss}ende Freigabe eines 64-Byte-Blocks aus dem Zephyr-Systemheap gemessen. Tabelle~\ref{tab:heap-laufzeiten} stellt die mittleren Laufzeiten beider Operationen gegen{\"u}ber.

\begin{longtable}{p{0.313\linewidth}p{0.313\linewidth}p{0.313\linewidth}}
\caption{Laufzeiten der Heap-Allokation und Heap-Freigabe}\label{tab:heap-laufzeiten}\\
\toprule
\textbf{Profil} & \textbf{Allokation} & \textbf{Freigabe} \\
\midrule
Minimal & 518,71 Zyklen & 312,82 Zyklen \\
Logging & 496,03 Zyklen & 319,02 Zyklen \\
IoT-Stack & 521,78 Zyklen & 318,99 Zyklen \\
\bottomrule
\end{longtable}

Tabelle~\ref{tab:heap-laufzeiten} zeigt, dass die Unterschiede der Allokation und Freigabe zwischen den Profilen klein und nicht monoton sind.

Die mittlere Allokationszeit liegt hierbei zwischen 496,03 und 521,78 Zyklen bzw. zwischen 7,75 und 8,15 \textmu{}s. Im Vergleich zu Minimal ist der Logging-Wert 4,37 \% niedriger, w{\"a}hrend der IoT-Wert 0,59 \% h{\"o}her liegt. Die Freigabe braucht zwischen 312,82 und 319,02 Zyklen bzw. zwischen 4,89 und 4,98 \textmu{}s. Logging und IoT liegen dabei jeweils ungef{\"a}hr 2 \% {\"u}ber dem Minimalprofil.

\section{Heap-Fragmentierung}
Der Fragmentierungstest benutzt einen eigenen Heap mit einer Gr{\"o}{\ss}e von 8192 Byte. Dieser wird zuerst mit 128-Byte-Bl{\"o}cken gef{\"u}llt. Anschlie{\ss}end wird jeder zweite erfolgreich angeforderte Block freigegeben. Tabelle~\ref{tab:heap-fragmentierung} dokumentiert die vier Zust{\"a}nde des Tests.

\begin{longtable}{p{0.235\linewidth}p{0.235\linewidth}p{0.235\linewidth}p{0.235\linewidth}}
\caption{Ergebnisse des Heap-Fragmentierungstests}\label{tab:heap-fragmentierung}\\
\toprule
\textbf{Zustand} & \textbf{Angefordert} & \textbf{Nicht angefordert} & \textbf{Gr{\"o}{\ss}te erfolgreiche Probe} \\
\midrule
Leer & 0 Byte & 8192 Byte & -- \\
Gef{\"u}llt & 7552 Byte & 640 Byte & 96 Byte \\
Fragmentiert & 3712 Byte & 4480 Byte & 128 Byte \\
Bereinigt & 0 Byte & 8192 Byte & -- \\
\bottomrule
\end{longtable}

Wie Tabelle~\ref{tab:heap-fragmentierung} zeigt, steigt der nicht angeforderte Speicher nach dem Fragmentieren deutlich an, ohne dass ein entsprechend gro{\ss}er zusammenh{\"a}ngender Block verf{\"u}gbar wird.

Nach dem Fragmentieren sind rechnerisch 4480 Byte nicht mehr angefordert. Dennoch kann nur ein zusammenh{\"a}ngender Block von 128 Byte erfolgreich angefordert werden. Die freien Bereiche sind daher {\"u}ber den Heap verteilt und k{\"o}nnen nicht zu einer gr{\"o}{\ss}eren zusammenh{\"a}ngenden Anforderung zusammengefasst werden. Alle drei Profile liefern gleiche Werte, weil der benutzte Heap nur dem Benchmark geh{\"o}rt und in jedem Profil dasselbe Allokationsmuster ausgef{\"u}hrt wird.

Die Angaben \glqq{}angefordert\grqq{} und \glqq{}nicht angefordert\grqq{} beschreiben dabei nur die Nutzbytes der Benchmark-Anforderungen. Interne Metadaten, Ausrichtung und Rundung des Allokators werden dadurch nicht erfasst.

\section{RAM-/ ROM-Bedarf}
Der statische Speicherbedarf wurde anhand der jeweiligen BENCH\_NONE-Referenzbauten bestimmt. Diese enthalten den gemeinsamen Messsockel, jedoch keinen bestimmten Benchmarkcode. Tabelle~\ref{tab:speicherbedarf-referenzbauten} vergleicht den Flash- und RAM-Bedarf.

\begin{longtable}{p{0.313\linewidth}p{0.313\linewidth}p{0.313\linewidth}}
\caption{Statischer Flash- und RAM-Bedarf der Referenzbauten}\label{tab:speicherbedarf-referenzbauten}\\
\toprule
Profil & Flash & RAM \\
\midrule
Minimal & 19.912 Byte & 69.592 Byte \\
Logging & 29.216 Byte & 71.712 Byte \\
IoT-Stack & 185.524 Byte & 141.040 Byte \\
\bottomrule
\end{longtable}

Tabelle~\ref{tab:speicherbedarf-referenzbauten} macht sichtbar, dass insbesondere das IoT-Profil den statischen Speicherbedarf stark erh{\"o}ht.

Das Logging-Profil braucht im Vergleich zu Minimal weitere 9304 Byte Flash und 2120 Byte RAM. Dies entspricht Zunahmen von 46,73 \% bzw. 3,05 \%. Beim IoT-Profil betr{\"a}gt der weitere Bedarf 165.612 Byte Flash und 71.448 Byte RAM. Im Vergleich zu Minimal steigt der Flash-Bedarf damit um 831,72 \% und der RAM-Bedarf um 102,67 \%.

Der RAM-Bedarf aller Profile enth{\"a}lt dabei den gemeinsam konfigurierten 64-KiB-Systemheap. Im IoT-Profil kommen unter anderem der 16-KiB-mbedTLS-Heap, Netzwerkpuffer, OpenThread-Strukturen und die vergr{\"o}{\ss}erten Main- und System-Workqueue-Stacks hinzu, die f{\"u}r eine funktionsf{\"a}hige Initialisierung gebraucht werden.
